\chapter{Algebras and Quaternions}

\section{Algebras}

\begin{definition}
  A {\it finite dimensional algebra \(A\)} is a finite dimesional vecotr space over \(\K\)
  which is also an associative and unial ring such that for all \(r,s\in\K\) and
  \(a,b\in A\), \((ra)(sb)=(rs)(ab)\).
\end{definition}

If every non-zero element \(u\in A\) is invertible, then \(A\) is a {\it division algebra over \(\K\)}.

\begin{definition}
  Let \(A\) and \(B\) be two algebras over \(\K\). A linear transformation \(\varphi:A\longrightarrow B\)
  that is also a ring homomorphism i scalled a {\it homomorphism of algebras}.
\end{definition}

\begin{definition}
  An algebra \(A\) is simple if it has only one proper two-sided ideal, namely
  (0).
\end{definition}

\begin{proposition}
  Let \(\K\) be a field.
  \begin{itemize}
    \item[\((i)\)] For a division algebra \(\mathbb{D}\) over \(\K\), \(\mathbb{D}\) is
      simple.
    \item[\((ii)\)] For a simple algebra \(A\) over \(\K\), \(M_n(A)\) is simple.
  \end{itemize}
\end{proposition}

\begin{theorem}[Wedderburn's Theorem]
  Let \(A\) be a finite dimensional simple algebra over \(\K\).
  \begin{itemize}
    \item[\((i)\)] \(A\) is isomorphic to some \(M_n(\mathbb{D})\).
    \item[\((ii)\)] If \(\K\) is algenraically closed, then \(A\) is isomorphic to some
      \(M_n(\K)\).
    \item[\((i)\)] Every finite dimensional division algebra \(\mathbb{D}\) over \(\K\) has
      dimension of the form \(\text{dim}_{\K}\mathbb{D}=d^2\) for some \(d\in\N\).
  \end{itemize}
\end{theorem}

\begin{proposition}
  A finite dimensional algebra \(A\) is a product of \(l\) algebras {\it iff} there are
  orthagonal central idempotents \(e_1,\dots,e_l\) such that \(e_1+\cdots+e_l=1\).
\end{proposition}

\begin{definition}
  An algebra \(A\) is {\it semi-simple} if it is a product of simple algebras \(A_k\),
  \[ A=A_1\times\cdots\times A_l \]
\end{definition}

\begin{proposition}
  A finite dimensional algebra \(A\) is semi-simple {\it iff} there are indecomposible orthagonal
  central idempotents \(e_1,\dots,e_l\) such that \(e_1+\cdots+e_l=1\).
\end{proposition}

\section{Linear Algebra over Division Algebras}
Let \(\D\) be a finite dimensional division algebra over a field \(\K\).

\begin{definition}
  A (right) vector spcae \(V\) over \(\D\) is an abelian group with a right scalar
  multiplication by elements of \(D\) so that for \(u,v\in V\), \(x,y\in\D\),
  \begin{eqnarray*}
    v(xy) &=& (vx)y \\
    v(x+y) &=& vx + vy \\
    (u+v)x &=& ux+vx \\
    v1 &=& v
  \end{eqnarray*}
\end{definition}

All obvious notions of linear transformations, subspace, kernels, and images make sense,
as do those of spanning sets and linear independence over \(\D\). The definitions of basis,
dimension, and matrices of linear transformatins are similar to that of vector spaces over
fields.

For a division algebra \(\D\), each matrix \(A\in M_n(\D)\) acts by multiplication on the
left of \(\D^n\). For any subfield \(\K_1\subseteq\D\) containing \(\K\), \(A\) induces
a (right) linear transformation over \(\K\), 
\[ \D^n\longrightarrow\D^n;\hspace{10px} x\mapsto Ax. \]
If we choose a basis over \(\K_1\) for \(\D\), \(A\) gives rise to a matrix \(\Lambda_A\in M_{nd}(\K_1)\)
where \(d=\text{dim}_{\K_1}\D\). \(A\) is invertible in \(M_n(\D)\) {\it iff} \(\Lambda_A\) is
invertible. Hence to determine invertibility of \(A\), it suffices to consider \(\text{det }\Lambda_A\).

The resulting function
\[ \text{Rdet}_{\K_1}: M_n(\D)\longrightarrow \K_1;\hspace{10px}\text{Rdet}_{\K_1}(A)=\text{det }\Lambda_A \]
is called the {\it reduced determinant} of \(M_n(\D)\), and it is a group homomorphism.

\section{Quaternions}

\begin{proposition}
  If \(A\) is a finite dimensional commutative division algebra over \(\R\),
  then either \(A=\R\) or there is an isomorphism of real algebra auch that \(A\cong\C\).
\end{proposition}

\begin{proof}
  Let \(\alpha\in A\). Since \(A\) is a finite dimensional real vector spaces, the powers
  of \(\alpha\) must be linearly dependent, sasy
  \[ t_0+t_1\alpha+\cdots+t_m\alpha^m=0 \]
  for some \(t_j\in\R\) with \(m\ge 1\) and \(t_m\ne 0\). We can choose \(m\) to be
  minimal with these properties. If \(t_0=0\), then
  \[ t_1+t_2\alpha+t_3\alpha^2+\cdots+t_m\alpha^{m-1}=0 \],
  contradicting the minimality of \(m\); so \(t_0\ne 0\). in fact, the polynomial
  \[ p(X)=t_0+t_1X+\cdots+t_mX^m \in\R[X] \]
  is irreducible. TO see this, wuppose that \(p(X)=p_1(X)p_2(X)\); then as \(A\) is a
  division algebra, either \(p_1(\alpha)=0\) or \(p_2(\alpha)=0\), contradicting minimality
  of \(m\) if both \(\text{deg }p_1(X)>0\) and \(\text{deg }p_2(X)>0\).

  Consider the real subspace
  \[ \R(\alpha)=\left\{\sum_{j=0}^k s_j\alpha^j:s_j\in\R\right\}\subseteq A. \]
  Then \(\R(\alpha)\) is easily seen to be a real subalgebra of \(A\) and the elements
  \(1,\alpha,\dots,\alpha^{m-1}\) form a basis, hence \(\text{dim}_{\R}\R(\alpha)=m\).

  Let \(\gamma\in\C\) be any complex root of the irreducible polynomial
  \[ t_0+t_1X+\cdots+t_mX^m \in\R[X]. \]
  There is a real linear transformation which is actualy an injection,
  \[ \varphi:\R(\alpha)\longrightarrow\C;\hspace{10px}\varphi\left(\sum_{j=0}^{m-1} s_j\alpha^j\right) = \sum_{j=0}^{m-1}s_j\gamma^j. \]
  It is easy to see thaat this is actually a homomorphism of real algebras. Hence
  \(\varphi(\R(\alpha))\subseteq\C\) is a subalgebra. But as \(\text{dim}_{\R}\C=2\),
  this implies that \(m=\text{dim}_{\R}\R(\alpha)\le 2\). If \(m=1\), \(\alpha\in\R\).
  If \(m=2\), then \(\varphi(\R(\alpha))=\C\).

  So either \(\text{dim}_{\R}A=1\) and \(A=\R\), or \(\text{dim}_{\R}A>1\) and we can
  choose an \(\alpha\in A\) with \(\C\cong\R(\alpha)\). This means that we can view \(A\)
  as a finite dimensional complex algebra. Now for any \(\beta\in A\) there is polynomial
  \[ q(X)=u_0+u_1X+\cdots+u_lX^l\in\C[X] \]
  with \(l\ge 1\) and \(u_l\ne 0\).Again choosing \(l\) to be minimal with this
  property, \(q(X)\) is irreducible. But then sinec \(q(X)\) has a root in \(\C\),
  \(l=1\) and \(\beta\in\C\). This shows that \(A=\C\) whenever \(\text{dim}_{\R}A>1\).
\end{proof}

However, the quesion of what finite dimesnional division algebra over \(\R\) exist is
less easy to decide. In fact, there is only one other upto an isomorphism, the skew-field
of quaternions, denoted \(\Q\). We will construct \(\Q\) as a ring of \(2\times 2\) complex
matrices.

Let
\[ \Q=\left\{ \left[\begin{array}{cc} z & w \\ -\bar{w} & \bar{z}\end{array}\right]:z,w\in\C \right\} \subseteq M_2(\C). \]
\(\Q\) is a subring of \(M_2(\C)\), and in fact a reaal subalgebra of \(M_2(\C)\).
It also contains a copy of \(\C\), namely the subalgebra
\[ \left\{\left[\begin{array}{cc} z & 0 \\ 0 & \bar{z} \end{array}\right]:z\in\C\right\}\subseteq\Q. \]
However, \(\Q\) is not an algebra over \(\C\).

An element of \(\Q\) is invertible {\it iff} it is an invertible as a matrix. It is useful
to denote a norm in \(\Q\) as
\[ |H|=\text{det }\left[\begin{array}{cc} z & w \\ -\bar{w} & \bar{z} \end{array}\right] = |z|^2+|w|^2. \]

Then \( SU(2)=\{A\in\Q:|A|=1\}\le\Q^{\times}\).

As a basis of \(\Q\) over \(\R\), we have the matrices
\[
  1=I,\hspace{10px}
  i=\left[\begin{array}{cc} i & 0 \\ 0 & -i \end{array}\right],\hspace{10px}
  j=\left[\begin{array}{cc} 0 & 1 \\ -1 & 0 \end{array}\right],\hspace{10px}
  k=\left[\begin{array}{cc} 0 & i \\ i & 0 \end{array}\right].
\]
These satisfy the equations
\[ i^2=j^2=k^2=-1,\hspace{10px} ij=k=-ji,\hspace{10px} jk=i=-kj,\hspace{10px} ki=j=-ik. \]

From now on we will write quaternions in the form
\[ q = xi + yj + zk +t1 \]
q is a pure quaternion {\it iff} t=0.

There is a quternionic analogue of complex conjugation, namely
\[ q = xi + yj + zk +t1 \mapsto  \bar{q} = -xi - yj - zk +t1\]

This satisfies the properties
\begin{eqnarray*}
  \overline{q_1+q_2} &=& \overline{q}_1+\overline{q}_2 \\
  \overline{q_1q_2} &=& \overline{q}_1\overline{q}_2 \\
  \overline{q}=q &\Longleftrightarrow& \text{ q is a real quaternion}; \\
  \overline{q}=-q &\Longleftrightarrow& \text{ q is a pure quaternion}.
\end{eqnarray*}

The inverse of a non-zero quaternion can be written as
\[ q^{-1}=\frac{\overline{q}}{q\overline{q}} \]

In terms of the matrix description of \(\Q\), quaternionic conjugation is given by the
hermitian conjugation,
\[
  \left[\begin{array}{cc}
      z & w \\
      -\bar{w} & \bar{z}
  \end{array}\right]
  \mapsto
  \left[\begin{array}{cc}
      \bar{z} & -w \\
      \bar{w} & z
  \end{array}\right].
\]

\section{Quaternionic Matrix Groups}
The above norm || on \(\Q\) extends to a norm on \(\Q^n\), viewed as a right \(\Q\)-vector
space. We can define a quaternionic inner product on \(\Q\) by
\[ x\cdot y=x^*y=\sum_{r=1}^n \bar{x}_ry_r, \]
where we define the {\it quaternionic conjugate} of a vector by
\[ \left[\begin{array}{c} x_1 \\ x_2 \\ \vdots \\ x_n \end{array}\right]^* = [\begin{array}{cccc} \bar{x}_1 & \bar{x}_2 & \cdots & \bar{x}_n \end{array}]. \]

Similarly, for any matrix \([a_{ij}]\) over \(\Q\) we can define \([a_{ij}]^*=[\bar{a}_{ji}].\)

The length of \(x\in\Q^n\) is defined to be \(|x|=\sqrt{x*x}\).

\section{Automorphism Groups of Algebras}

\[ \blacktriangle \blacktriangledown \blacktriangle \]
